% Source: https://artofproblemsolving.com/community/c3h2849536_arml_relay_problems
% Stable identity: aops-topic-2849536-q03
The vertices lie on a circle, with consecutive vertices $6^\circ$ apart. The chords $A_1A_{20}$ and $A_{10}A_{59}$ cross inside the circle (the endpoints alternate around the circle: $A_{59},A_1,A_{10},A_{20}$).
\begin{asy}
usepackage("amsmath");
size(4.5cm);
defaultpen(fontsize(10pt));
pair P(int k){return dir(90+6*(k-1));}
draw(circle((0,0),1),gray);
draw(P(1)--P(20),linewidth(0.9));draw(P(10)--P(59),linewidth(0.9));
pair X=extension(P(1),P(20),P(10),P(59));dot(X);
dot(P(1));dot(P(20));dot(P(10));dot(P(59));
label("$A_1$",P(1),N);label("$A_{20}$",P(20),W);label("$A_{10}$",P(10),NW);label("$A_{59}$",P(59),NE);
\end{asy}

The angle between two chords crossing inside a circle is half the sum of the intercepted arcs. The arc $A_1A_{10}$ is $9\cdot6^\circ=54^\circ$ and the arc $A_{20}A_{59}$ (the other side) is $39\cdot6^\circ=234^\circ$. So one angle is $\dfrac{54^\circ+234^\circ}2=144^\circ$, and the acute angle is $180^\circ-144^\circ=\boxed{36^\circ}$.
